\section{Onde aprofundar}

\subsection*{Livros e artigos}
\small
\begin{longtable}{@{}p{4.4cm} p{3cm} p{4.8cm}@{}}
\toprule
Fonte & Autor(es), ano & Por que ler \\
\midrule \endhead
\citetitle{livro:martin2017} \pt{} (\emph{Arquitetura Limpa})
  & \citeauthor{livro:martin2017}, \citeyear{livro:martin2017}
  & Caps.~22 (os círculos e a Dependency Rule) e 32 (o framework como
  detalhe) \\
\citetitle{artigo:cockburn-hexagonal}
  & \citeauthor{artigo:cockburn-hexagonal}, 2005
  & O texto original de ports \& adapters, com a motivação de testar a
  aplicação sem UI nem banco \\
\bottomrule
\end{longtable}
\normalsize

\subsection*{Documentação}
\citetitle{doc:quarkus-cdi}, \citetitle{doc:quarkus-cdi-reference},
\citetitle{doc:quarkus-rest}, \citetitle{doc:quarkus-panache},
\citetitle{doc:quarkus-config-mappings},
\citetitle{doc:quarkus-testing}, \citetitle{doc:quarkus-dev-services},
\citetitle{doc:quarkus-fault-tolerance}, \citetitle{doc:quarkus-health},
\citetitle{doc:quarkus-micrometer} e
\citetitle{doc:quarkus-opentelemetry}.

\subsection*{No IDE: blog da JetBrains}
Outros posts da JetBrains no apêndice~\ref{ap:jetbrains}.
\begin{itemize}
\obra*{artigo:jetbrains-testcontainers-spring-boot}{Testes de integração
  com Testcontainers; o exemplo é Spring Boot, mas a ideia é a mesma dos Dev
  Services do Quarkus.}
\obra*{curso:jetbrains-guide-javax-jakarta}{Migrar de \texttt{javax} para
  \texttt{jakarta}, o namespace que o Quarkus usa.}
\end{itemize}

\subsection*{Pesquisa acadêmica}
Não encontramos estudo acadêmico revisado por pares sobre arquitetura
hexagonal ou Clean Architecture; os itens abaixo são a base conceitual e
empírica por trás delas. Acesso e links no apêndice~\ref{ap:pesquisa}.
\begin{itemize}
\pesquisa{paper:parnas-modules-1972}{A raiz de portas e adaptadores:
  esconder atrás de uma interface o que pode mudar.}
\pesquisa{paper:garlan-shaw-intro-softarch-1994}{Relatório técnico da CMU
  sobre estilos arquiteturais: camadas, eventos, pipes and filters.}
\pesquisa{paper:perry-wolf-foundations-1992}{Define arquitetura e introduz
  as ideias de erosão e desvio arquitetural.}
\pesquisa{paper:desilva-erosion-survey-2012}{Técnicas contra a erosão, como
  regras de dependência no estilo do ArchUnit.}
\pesquisa{paper:dragoni-microservices-2017}{Panorama histórico de
  microsserviços e dos problemas em aberto.}
\pesquisa{paper:soldani-microservices-pains-gains-2018}{O que praticantes
  relatam de bom e de ruim em microsserviços.}
\pesquisa{paper:viggiato-microservices-practice-2018}{Pesquisa da UFMG com
  122 profissionais sobre microsserviços na prática.}
\pesquisa{paper:spadini-mock-objects-java-2019}{Por que se mocka em Java e
  como os mocks evoluem com o código.}
\pesquisa{paper:almeida-mock-objects-sbqs-2023}{Da UFF: Mockito é o
  framework mais usado em Java e dependências externas são as mais mockadas.}
\pesquisa*{livro:valente2020}{Cap.~7 (arquitetura: camadas, MVC,
  microsserviços, mensageria) e cap.~8 (testes).}
\end{itemize}
